%% =========================================================
%% Chapter: Physics examples of the multi-block asymmetric
%% compression theorem (continued).
%% Covers the anomalous Z_3 symmetry and its condensation defect
%% (Example I of Notes/OpenProblemsTN/checks/asym_z3_anomalous_data.md).
%% =========================================================

\section{The anomalous \texorpdfstring{$\Z_3$}{Z3} symmetry and its condensation defect}%
\label{sec:asymex_z3anomalous}

On a periodic chain of $L>0$ qutrits, write $X\ket{s}=\ket{s+1}$ for the cyclic shift and
$\omega=e^{2\pi i/3}$. The phase-decorated shift acts on a basis vector by
\begin{align}
    U\ket{s_1,\ldots,s_L}
    & =\left(\prod_{k=1}^{L}\omega^{[s_k\neq0]\,s_{k+1}}\right)
    X^{\otimes L}\ket{s_1,\ldots,s_L},
    \label{eq:asymex_z3anomalous_operator}
\end{align}
where $s_{L+1}=s_1$, addition is modulo three, and $[s_k\neq0]$ is the indicator of
$s_k\neq0$. This operator has a matrix product representation of bond dimension two and
satisfies $U^3=\id$ at every positive length. Together with the identity and its inverse,
it therefore gives a representation of $\Z_3$. The anomaly class belongs to
$H^3(\Z_3,U(1))=\Z_3$ and is determined by the associator of the fusion tensors, as in
Garre-Rubio and Schuch~\cite{GarreRubioSchuch2024DomainWall}. Its computation is separate
from the compression and intertwiner calculations below. These calculations give the
group law, the absence of nonzero sitewise intertwiners for the four products of
nonidentity tensors, and the periodic identity $O_L(A)^2=3O_L(A)$ for the condensation
defect $A=\delta\oplus M\oplus M^\dagger$.

\begin{definition}[The tensors of the representation]
    \label{def:asymex_z3anomalous_tensors}
    \lean{Z3Anomalous.uTensor, Z3Anomalous.uDagTensor, Z3Anomalous.identityTensor,
        Z3Anomalous.uMPS, Z3Anomalous.uDagMPS, Z3Anomalous.identityMPS}
    \leanok
    \uses{def:mpo_tensor, def:mpo_to_mps, def:mpo_family}
    The tensor $M$ of $U$ has bond dimension two, with
    \begin{align}
        M^{1,0} & =\begin{pmatrix}1&0\\1&0\end{pmatrix},        &
        M^{2,1} & =\begin{pmatrix}0&1\\0&\omega\end{pmatrix},   &
        M^{0,2} & =\begin{pmatrix}0&1\\0&\omega^2\end{pmatrix},
        \notag
    \end{align}
    and all other matrices zero. The tensor denoted by $M^\dagger$ has bond dimension two,
    with
    \begin{align}
        (M^\dagger)^{0,1} & =\begin{pmatrix}1&0\\1&0\end{pmatrix},        &
        (M^\dagger)^{1,2} & =\begin{pmatrix}0&1\\0&\omega^2\end{pmatrix}, &
        (M^\dagger)^{2,0} & =\begin{pmatrix}0&1\\0&\omega\end{pmatrix},
        \notag
    \end{align}
    and all other matrices zero. The identity tensor $\delta$ has scalar matrices
    $\delta^{ij}$. Read over the alphabet of index pairs, $M$, $M^\dagger$ and $\delta$
    are matrix product state tensors of physical dimension nine, with matrices $M^{ij}$,
    $(M^\dagger)^{ij}$, $\delta^{ij}$ at the letter $3i+j$.
    Diagrammatically, $U=O_L(M)$ is the periodic operator
    \begin{center}
        % Formula: U = X^{\otimes L} \prod_k \omega^{[s_k \neq 0] s_{k+1}}, the periodic
        % operator O_L(M) generated by the tensor M above.
        % Ink-to-index map: each node is M; its upper leg carries the ket index s_k+1 and its
        % lower leg the bra index s_k; the horizontal legs are the bond-two virtual index.
        % Contraction set: every horizontal virtual bond is contracted, including the return
        % bond that closes the ring.
        % Boundary signature: virtual-west=0 | virtual-east=0 | physical-up=L |
        % physical-down=L; three representative sites are drawn, the remaining L-3 abbreviated
        % by the ellipsis.
        \tenkzkernel
        \begin{tenkz}[rows={wire}, cols=4, boundary=periodic]
            \tn[skin=mpo, ports={90:physical:$s_1{+}1$, 270:physical:$s_1$}]{M} &
            \tn[skin=mpo, ports={90:physical:$s_2{+}1$, 270:physical:$s_2$}]{M} &
            \tn[skin=dots]{} &
            \tn[skin=mpo, ports={90:physical:$s_L{+}1$, 270:physical:$s_L$}]{M}
        \end{tenkz}
    \end{center}
    generated by $M$.
\end{definition}

\begin{remark}[Notation for the group inverse]
    \label{rem:asymex_z3anomalous_notation}
    The symbol $M^\dagger$ denotes the independently specified tensor in
    Definition~\ref{def:asymex_z3anomalous_tensors}, not the conjugate transpose of each
    bond matrix of $M$. The results below establish
    $O_L(M)O_L(M^\dagger)=O_L(M^\dagger)O_L(M)=\id$ and $O_L(M)^3=\id$ for $L>0$.
    Thus $O_L(M^\dagger)=U^{-1}=U^2$. Identifying this inverse with the operator adjoint
    also uses the unitarity of the phase-decorated shift.
\end{remark}

\begin{theorem}[Normality of the three blocks up to normalization]
    \label{thm:asymex_z3anomalous_normal}
    \lean{Z3Anomalous.uMPS_isNormal, Z3Anomalous.uDagMPS_isNormal,
        Z3Anomalous.identityMPS_isNormal}
    \leanok
    \uses{def:asymex_z3anomalous_tensors, def:normal}
    The length-two words of $M$ span $\MN{2}$, as do those of $M^\dagger$.
    The length-one words of $\delta$ span $\MN{1}$. Thus each tensor is proportional
    to a normal tensor (NT); the NT convention additionally normalizes the spectral
    radius of the associated completely positive map to one.
\end{theorem}

\begin{proof}
    \leanok
    For $M$, four length-two words in the two letters $(0,2)$ and $(1,0)$ combine, with
    coefficients in the Eisenstein integers $\Z[\omega]$, to $(\omega-1)$ times each matrix unit
    of $\MN{2}$; since $\omega-1\neq0$, the length-two words already span $\MN{2}$, which is a
    finite verification over $\Z[\omega]$. The same holds for $M^\dagger$ with the letters
    $(0,1)$ and $(1,2)$. For the identity tensor, a diagonal physical index pair gives the
    nonzero scalar $1$.
\end{proof}

\section{Fusion of the anomalous representation}%
\label{sec:asymex_z3anomalous_fusion}

The stacked product of two matrix product operator tensors $M_g$, $M_h$ is
$B^{ij}=\sum_kM_g^{ik}\otimes M_h^{kj}$, of bond dimension the product of the two bond
dimensions (Theorem~\ref{thm:asym_mpo_product}). Write $M_0=\delta$, $M_1=M$ and
$M_2=M^\dagger$, with group indices added modulo three. For the representation of
Definition~\ref{def:asymex_z3anomalous_tensors}, the periodic operator of the stacked
product $M_g\cdot M_h$ equals $O_L(M_{g+h})$ for $L>0$. In particular, stacking $M$ over
$M$ gives $U^2=O_L(M^\dagger)$:
\begin{center}
    % Formula: O_L(U)O_L(U) = O_L(U^\dagger) (Theorem~\ref{thm:asymex_z3anomalous_uu_trace}),
    % the stacked product B=M\cdot M contracted with the periodic-operator identity.
    % Ink-to-index map: on the left, the upper row is M with ket index i_k and bra index m_k,
    % the lower row is M with ket index m_k and bra index j_k; the shared leg m_k is contracted.
    % On the right, the single row is M^\dagger with ket index i_k and bra index j_k.
    % Contraction set: every vertical physical pair m_k, every horizontal virtual bond of each
    % row, and the periodic return bonds of both rings.
    % Boundary signature: both sides (0,0,L,L); the drawn three-site representative emits
    % (0,0,3,3).
    \tenkzkernel
    \begin{tenkz}[rows={op,op}, cols=4, boundary=periodic]
        \tn[skin=mpo, ports={90:physical:$i_1$, 270:physical}]{M} &
        \tn[skin=mpo, ports={90:physical:$i_2$, 270:physical}]{M} &
        \tn[skin=dots, ports={180:virtual, 0:virtual}]{} &
        \tn[skin=mpo, ports={90:physical:$i_L$, 270:physical}]{M} \\
        \tn[skin=mpo, ports={90:physical, 270:physical:$j_1$}]{M} &
        \tn[skin=mpo, ports={90:physical, 270:physical:$j_2$}]{M} &
        \tn[skin=dots, ports={180:virtual, 0:virtual}]{} &
        \tn[skin=mpo, ports={90:physical, 270:physical:$j_L$}]{M}
    \end{tenkz}
    \quad$=$\quad
    \begin{tenkz}[rows={wire}, cols=4, boundary=periodic]
        \tn[skin=mpo, ports={90:physical:$i_1$, 270:physical:$j_1$}]{M^\dagger} &
        \tn[skin=mpo, ports={90:physical:$i_2$, 270:physical:$j_2$}]{M^\dagger} &
        \tn[skin=dots]{} &
        \tn[skin=mpo, ports={90:physical:$i_L$, 270:physical:$j_L$}]{M^\dagger}
    \end{tenkz}
\end{center}
The other nontrivial products have targets $\delta$ for $M\cdot M^\dagger$ and
$M^\dagger\cdot M$, and $M$ for $M^\dagger\cdot M^\dagger$.

\begin{theorem}[Compression of \texorpdfstring{$M\cdot M$ and $M^\dagger\cdot M^\dagger$}{the two equal-factor products}]
    \label{thm:asymex_z3anomalous_uu_compression}
    \lean{Z3Anomalous.uu_compression, Z3Anomalous.uu_isReduction, Z3Anomalous.uu_right_eq,
        Z3Anomalous.uu_left_eq, Z3Anomalous.uu_dim_eq,
        Z3Anomalous.uu_evalWord_remainder_eq_zero, Z3Anomalous.uu_remainder_mul_ne_zero,
        Z3Anomalous.dd_compression, Z3Anomalous.dd_isReduction, Z3Anomalous.dd_right_eq,
        Z3Anomalous.dd_left_eq, Z3Anomalous.dd_dim_eq,
        Z3Anomalous.dd_evalWord_remainder_eq_zero, Z3Anomalous.dd_remainder_mul_ne_zero}
    \leanok
    \uses{def:asymex_z3anomalous_tensors, thm:asym_multiblock}
    The stacked tensor $B=M\cdot M$ has the block triangular form of
    Theorem~\ref{thm:asym_multiblock} for the single target $M^\dagger$ with $z=2$
    one-dimensional zero diagonal blocks, $4=2+2$. The compression pair is
    \begin{align}
        V & =\begin{pmatrix}0&0\\0&-\omega\\0&1+\omega\\-1&0\end{pmatrix}, &
        W & =\begin{pmatrix}0&0&0&-1\\0&0&-\omega&0\end{pmatrix},
        \notag
    \end{align}
    with $WV=\id$ and $WB^wV=(M^\dagger)^w$ for every word, and the remainder
    $R^{ij}=B^{ij}-V(M^\dagger)^{ij}W$ is nilpotent of order exactly three: some product of two
    of its letters is nonzero, but every product of three vanishes. The stacked tensor
    $M^\dagger\cdot M^\dagger$ has the analogous block triangular form for the single
    target $M$, again with $z=2$ one-dimensional zero diagonal blocks and a remainder
    of order exactly three.
\end{theorem}

\begin{proof}
    \leanok
    \uses{thm:asymex_explicit_gauge}
    Conjugating the letters of $B$ by the change of bond coordinates whose columns are, in the
    order of the flag, the common kernel of every letter, an adapted basis of the common image
    modulo that kernel, and a completion, makes every letter block upper triangular with
    $M^\dagger$ on the matched diagonal block; this is a finite verification over $\Z[\omega]$,
    transported to the complex numbers along the entrywise embedding, exactly as in
    Theorem~\ref{thm:asymex_explicit_gauge}. The compression pair and the sharp nilpotency order
    follow by reading off the corresponding rows and columns of the gauge and checking the two
    remaining finite identities. The block for $M^\dagger\cdot M^\dagger$ is the same
    computation with the roles of $M$ and $M^\dagger$ exchanged.
\end{proof}

\begin{theorem}[Periodic fusion \texorpdfstring{$U\cdot U=U^\dagger$ and $U^\dagger\cdot U^\dagger=U$}{of the nonidentity group elements}]
    \label{thm:asymex_z3anomalous_uu_trace}
    \lean{Z3Anomalous.uu_trace_evalWord, Z3Anomalous.dd_trace_evalWord, Z3Anomalous.mpo_uu,
        Z3Anomalous.mpo_dd}
    \leanok
    \uses{thm:asymex_z3anomalous_uu_compression, prop:asym_compression_trace,
        thm:asym_mpo_product}
    For every nonempty word $w$, $\tr(B^w)=\tr((M^\dagger)^w)$, where $B=M\cdot M$.
    Similarly, $\tr((M^\dagger\cdot M^\dagger)^w)=\tr(M^w)$. Hence
    $O_L(M)O_L(M)=O_L(M^\dagger)$ and
    $O_L(M^\dagger)O_L(M^\dagger)=O_L(M)$ at every positive length $L$.
\end{theorem}

\begin{proof}
    \leanok
    Proposition~\ref{prop:asym_compression_trace} applied to
    Theorem~\ref{thm:asymex_z3anomalous_uu_compression} gives the trace identities, and
    Theorem~\ref{thm:asym_mpo_product} gives the stated identities of periodic operators.
\end{proof}

\begin{theorem}[Compression of \texorpdfstring{$M\cdot M^\dagger$ and $M^\dagger\cdot M$}{the inverse-factor products}]
    \label{thm:asymex_z3anomalous_ud_compression}
    \lean{Z3Anomalous.ud_compression, Z3Anomalous.ud_isReduction, Z3Anomalous.ud_right_eq,
        Z3Anomalous.ud_left_eq, Z3Anomalous.ud_dim_eq,
        Z3Anomalous.ud_evalWord_remainder_eq_zero, Z3Anomalous.ud_remainder_mul_ne_zero,
        Z3Anomalous.du_compression, Z3Anomalous.du_isReduction, Z3Anomalous.du_right_eq,
        Z3Anomalous.du_left_eq, Z3Anomalous.du_dim_eq,
        Z3Anomalous.du_evalWord_remainder_eq_zero, Z3Anomalous.du_remainder_mul_ne_zero}
    \leanok
    \uses{def:asymex_z3anomalous_tensors, thm:asym_multiblock}
    The stacked tensor $B=M\cdot M^\dagger$ has the block triangular form of
    Theorem~\ref{thm:asym_multiblock} for the single target $\delta$ with $z=3$
    one-dimensional zero diagonal blocks, $4=1+3$. The compression pair satisfies
    $WV=\id$ and $WB^wV=\delta^w$ for every word, and the remainder is nilpotent of
    order exactly three, one short of the general bound $|S|+z=4$. The stacked tensor
    $M^\dagger\cdot M$ has the analogous block triangular form for the same target $\delta$.
\end{theorem}

\begin{proof}
    \leanok
    \uses{thm:asymex_explicit_gauge}
    Use the same conjugation argument as in
    Theorem~\ref{thm:asymex_z3anomalous_uu_compression}, over $\Z[\omega]$, with the flag
    now containing two one-dimensional zero diagonal blocks before the target and one after.
\end{proof}

\begin{theorem}[The exact \texorpdfstring{$\Z_3$}{Z3} representation]
    \label{thm:asymex_z3anomalous_group_law}
    \lean{Z3Anomalous.ud_trace_evalWord, Z3Anomalous.du_trace_evalWord, Z3Anomalous.mpo_ud,
        Z3Anomalous.mpo_du, Z3Anomalous.mpo_u_mul_uDag, Z3Anomalous.mpo_uDag_mul_u,
        Z3Anomalous.mpo_u_pow_three, Z3Anomalous.mpo_identityTensor}
    \leanok
    \uses{thm:asymex_z3anomalous_ud_compression, thm:asymex_z3anomalous_uu_trace,
        prop:asym_compression_trace, thm:asym_mpo_product}
    The periodic operator of the identity tensor is the identity map at every length,
    $O_L(\delta)=\id$. For every nonempty word $w$, both
    $\tr((M\cdot M^\dagger)^w)$ and $\tr((M^\dagger\cdot M)^w)$ equal $\tr(\delta^w)$.
    Thus $O_L(M)O_L(M^\dagger)=O_L(M^\dagger)O_L(M)=\id$ at every positive length.
    Together with $O_L(M)O_L(M)=O_L(M^\dagger)$, this gives $O_L(M)^3=\id$:
    the symmetry gives a representation of $\Z_3$ without any length-dependent phase.
\end{theorem}

\begin{proof}
    \leanok
    $O_L(M)O_L(M^\dagger)=\id$ and $O_L(M^\dagger)O_L(M)=\id$ follow as in the proof of
    Theorem~\ref{thm:asymex_z3anomalous_uu_trace}, from
    Theorem~\ref{thm:asymex_z3anomalous_ud_compression} together with $O_L(\delta)=\id$. Then
    $O_L(M)^3=O_L(M)^2O_L(M)=O_L(M^\dagger)O_L(M)=\id$.
\end{proof}

\begin{theorem}[Both sitewise intertwiner spaces vanish on every nontrivial block]
    \label{thm:asymex_z3anomalous_no_intertwiner}
    \lean{Z3Anomalous.uu_right_intertwiner_eq_zero, Z3Anomalous.uu_left_intertwiner_eq_zero,
        Z3Anomalous.dd_right_intertwiner_eq_zero, Z3Anomalous.dd_left_intertwiner_eq_zero,
        Z3Anomalous.ud_right_intertwiner_eq_zero, Z3Anomalous.ud_left_intertwiner_eq_zero,
        Z3Anomalous.du_right_intertwiner_eq_zero, Z3Anomalous.du_left_intertwiner_eq_zero}
    \leanok
    \uses{def:asymex_z3anomalous_tensors}
    For each of the four stacked products $M\cdot M$, $M\cdot M^\dagger$,
    $M^\dagger\cdot M$, $M^\dagger\cdot M^\dagger$, let $B$ denote the stacked tensor
    and let $C$ denote its respective target $M^\dagger$, $\delta$, $\delta$, $M$.
    There is no nonzero $X$ with $B^iX=XC^i$ for every letter $i$, and no nonzero $Y$
    with $YB^i=C^iY$ for every letter $i$. In particular, the compression identities
    $WB^wV=C^w$ do not give sitewise intertwiners.
\end{theorem}

\begin{proof}
    \leanok
    For each block, a selected set of eight sitewise equations has a coefficient matrix over
    $\Z[\omega]$ with an explicit left inverse up to a nonzero scalar; this forces every
    coordinate of $X$, respectively $Y$, to vanish. Each inverse is a finite verification over
    $\Z[\omega]$.
\end{proof}

\begin{remark}[Failure of closure under conjugate transposition]
    \label{rem:asymex_z3anomalous_anomaly}
    \uses{cor:asym_anomaly_not_star}
    By Corollary~\ref{cor:asym_anomaly_not_star}, Theorem~\ref{thm:asymex_z3anomalous_no_intertwiner}
    shows that the algebra generated by the matrices of each of the four stacked products is not
    closed under conjugate transposition. This conclusion does not calculate the anomaly
    class. Determining that class requires a separate computation of the $3$-cocycle
    associator of the fusion tensors, as in Garre-Rubio and Schuch
    \cite{GarreRubioSchuch2024DomainWall}; neither the absence of sitewise intertwiners nor
    the nilpotent remainder alone determines an element of $H^3(\Z_3,U(1))$.
\end{remark}

\section{The condensation defect}%
\label{sec:asymex_z3anomalous_defect}

\begin{definition}[The condensation defect]
    \label{def:asymex_z3anomalous_defect}
    \lean{Z3Anomalous.summandDim, Z3Anomalous.summandMPS, Z3Anomalous.defectTensor,
        Z3Anomalous.defectMPS}
    \leanok
    \uses{def:asymex_z3anomalous_tensors, def:mpdo_operator_product}
    The condensation defect is the block-diagonal tensor $A=\delta\oplus M\oplus M^\dagger$
    of bond dimension five, with the identity tensor, the tensor of $U$ and the tensor of
    $U^\dagger$ on its diagonal, in that order.
\end{definition}

\begin{theorem}[The defect is the sum of the three group elements]
    \label{thm:asymex_z3anomalous_defect_sum}
    \lean{Z3Anomalous.defectMPS_trace_evalWord, Z3Anomalous.mpo_defectTensor}
    \leanok
    \uses{def:asymex_z3anomalous_defect}
    For every word $w$, $\tr(A^w)=\tr(\delta^w)+\tr(M^w)+\tr((M^\dagger)^w)$; that is,
    $O_L(A)=O_L(\delta)+O_L(M)+O_L(M^\dagger)$ at every length $L$.
\end{theorem}

\begin{proof}
    \leanok
    The trace of a block-diagonal matrix is the sum of the traces of its diagonal blocks, and
    $A$ is block diagonal with the three summands as its blocks at every letter.
\end{proof}

\begin{proposition}[Block decomposition of the stacked square]
    \label{prop:asymex_z3anomalous_defect_square_blocks}
    \lean{Z3Anomalous.defectSquare, Z3Anomalous.defectSquareEis_eq, Z3Anomalous.pairStackEis,
        Z3Anomalous.squareBond}
    \leanok
    \uses{def:asymex_z3anomalous_defect, def:mpdo_operator_product}
    Write $M_0=\delta$, $M_1=M$, $M_2=M^\dagger$, and $B_{gh}=(M_g\cdot M_h)^a$ for the
    pair block at a fixed letter $a$ of the pair alphabet. After reindexing the bond
    coordinate $p=5p_1+p_2$ by the summand pair and its internal coordinates, the stacked
    tensor $A\cdot A$ of bond dimension twenty-five is block diagonal over the nine pairs
    $(g,h)\in\{0,1,2\}^2$, with $(g,h)$ block $M_g\cdot M_h$. In these coordinates,
    for every letter $a$,
    \begin{align}
        (A\cdot A)^{a} & =
        \begin{pmatrix}
            B_{00}                                                                   \\
             & B_{01}                                                                \\
             &        & B_{02}                                                       \\
             &        &        & B_{10}                                              \\
             &        &        &        & B_{11}                                     \\
             &        &        &        &        & B_{12}                            \\
             &        &        &        &        &        & B_{20}                   \\
             &        &        &        &        &        &        & B_{21}          \\
             &        &        &        &        &        &        &        & B_{22}
        \end{pmatrix},
        \label{eq:asymex_z3anomalous_defect_square_blocks}
    \end{align}
    with all off-diagonal blocks zero.
\end{proposition}

\begin{proof}
    \leanok
    Both sides are the same reindexing, along the bijection labelling the twenty-five bond
    coordinates by the nine pairs of coordinates of the summands, of the same finite family of
    matrices; this is a finite verification over $\Z[\omega]$.
\end{proof}

\begin{theorem}[The defect squares to three times itself]
    \label{thm:asymex_z3anomalous_defect_square}
    \lean{Z3Anomalous.defectSquare_trace_evalWord, Z3Anomalous.mpo_defect_mul_defect}
    \leanok
    \uses{prop:asymex_z3anomalous_defect_square_blocks, thm:asymex_z3anomalous_defect_sum,
        prop:asym_compression_trace}
    For every nonempty word $w$, $\tr((A\cdot A)^w)=3\bigl(\tr(\delta^w)+\tr(M^w)+
        \tr((M^\dagger)^w)\bigr)$; that is, $O_L(A)O_L(A)=3\,O_L(A)$ at every positive length $L$.
\end{theorem}

\begin{proof}
    \leanok
    By Proposition~\ref{prop:asymex_z3anomalous_defect_square_blocks}, the trace of
    $(A\cdot A)^w$ is the sum, over the nine pairs $(g,h)$, of the trace of $(M_g\cdot M_h)^w$;
    each of the nine summands equals the trace of $M_{g+h}^w$, by
    Theorem~\ref{thm:asymex_z3anomalous_uu_trace} and
    Theorem~\ref{thm:asymex_z3anomalous_group_law} on the four nontrivial pairs and trivially on
    the five pairs with $g=0$ or $h=0$, and $\delta$, $M$, $M^\dagger$ each occur as $M_{g+h}$
    for exactly three of the nine pairs. Summing gives the stated trace identity, and the
    operator-level argument of Theorem~\ref{thm:asymex_z3anomalous_uu_trace} gives the periodic
    operator identity.
\end{proof}

\begin{theorem}[Compression of the stacked square]
    \label{thm:asymex_z3anomalous_defect_compression}
    \lean{Z3Anomalous.pairTarget, Z3Anomalous.defectSquare_compression,
        Z3Anomalous.defectSquare_isReduction, Z3Anomalous.defectSquare_left_mul_right_of_ne,
        Z3Anomalous.defectSquare_dim_eq, Z3Anomalous.defectSquare_evalWord_remainder_eq_zero,
        Z3Anomalous.sum_pairTargetDim}
    \leanok
    \uses{prop:asymex_z3anomalous_defect_square_blocks, thm:asym_multiblock,
        thm:asymex_z3anomalous_uu_compression, thm:asymex_z3anomalous_ud_compression,
        thm:asymex_explicit_gauge}
    The stacked tensor $A\cdot A$ has the block triangular form of
    Theorem~\ref{thm:asym_multiblock} for the nine targets $M_{g+h}$, $(g,h)\in\{0,1,2\}^2$,
    with $z=10$ one-dimensional zero diagonal blocks and $25=15+10$. The nine compression
    pairs are mutually biorthogonal, and the remainder is nilpotent of order at most nineteen.
\end{theorem}

\begin{proof}
    \leanok
    \uses{thm:asymex_explicit_gauge}
    On the five pairs with $g=0$ or $h=0$, the pair block already equals the target with no
    zero diagonal block. On the four remaining pairs, apply the gauge of
    Theorem~\ref{thm:asymex_z3anomalous_uu_compression} or
    Theorem~\ref{thm:asymex_z3anomalous_ud_compression}. The direct sum of these five identity
    gauges and four block gauges, transported along the labelling of
    Proposition~\ref{prop:asymex_z3anomalous_defect_square_blocks}, is the gauge required by
    Theorem~\ref{thm:asym_multiblock}; the dimension count adds
    $1{+}2{+}2{+}2{+}4{+}4{+}2{+}4{+}4=25$ against target sizes
    $1{+}2{+}2{+}2{+}2{+}1{+}2{+}1{+}2=15$ and
    $0{+}0{+}0{+}0{+}2{+}3{+}0{+}3{+}2=10$ one-dimensional zero diagonal blocks.
    Clause (vi) of Theorem~\ref{thm:asym_multiblock} bounds the nilpotency order by
    $|S|+z=9+10=19$.
\end{proof}
